% © 2026 Ing. David Jaroš — CC BY-NC-ND 4.0
%
% This work is licensed under a Creative Commons Attribution-NonCommercial-NoDerivatives
% 4.0 International License (CC BY-NC-ND 4.0).
%
% License History: Earlier drafts (up to v0.3) were released under CC BY 4.0.
% From v0.4 onward, all material is released under CC BY-NC-ND 4.0 to protect
% the integrity of the theoretical work during ongoing academic development.
%
% See LICENSE.md for full license text.

\ifdefined\INCLUDEMODE
\else
  \documentclass[12pt,a4paper]{article}
  \usepackage[a4paper, margin=2.5cm]{geometry}
  \usepackage{amsmath,amssymb}
  \usepackage{tcolorbox}
  \newtcolorbox{statusbox}[1][]{colback=gray!8!white,colframe=black!60,
    arc=3pt,boxrule=1pt,fonttitle=\bfseries,title=Status Summary,#1}
  \title{Synthesis: Why $p=137$ and not $p=139$}
  \author{Ing.\ David Jaro\v{s}}
  \date{2026-05-10}
  \begin{document}
  \maketitle
\fi

\section{Why $p=137$ and not $p=139$}

Among primes with $g(X_0(p))=11$, i.e. $p\in\{131,137,139\}$:

\begin{enumerate}
\item \textbf{Congruence condition (Kronecker symbol).}
$\nu_2(p)>0 \iff p\equiv1\pmod4$ is proved arithmetically [L0].
This keeps only $p=137$ inside the genus-11 cluster.
Physical lift from UBT charge-conjugation $C$ to Atkin--Lehner fixed points
remains [OPEN].

\item \textbf{Self-consistency sign condition.}
For
$f(p)=\frac{5p-1}{p+1}-\ln p$,
\begin{equation}
f(137)=+0.036541,\qquad f(139)=+0.022669.
\end{equation}
Both are positive, so the sign test does \emph{not} separate 137 from 139.
Status: FAILED as selection criterion [L0 numerical check].

\item \textbf{QED residuum identification.}
$f(137)$ is close to observed $\Delta\alpha^{-1}$, but one-loop leptonic
$\sum\Delta_{\mathrm{QED}}(\Lambda=M_{\mathrm{Pl}})=0.114595$ is far from $f(137)$.
Status: [NC] in current implementation.

\item \textbf{Hecke prime specificity.}
$a_{137}(N=76,k=2)=-11$ is known, but $a_{131},a_{139}$ for the same form are
currently unavailable in this environment (LMFDB DNS blocked, no Sage/PARI).
Status: [OPEN]/UNKNOWN.
\end{enumerate}

Criteria that currently yield $p=137$ uniquely:
\begin{itemize}
  \item Arithmetic congruence filter ($g=11$ and $p\equiv1\pmod4$),
  \item with physical UBT justification of the $C\leftrightarrow w_p$ bridge still open.
\end{itemize}

\begin{statusbox}
Počet algebraických kritérií selektujících $p=137$: 1 z 4 (strictly proved selection),
plus 1 conditional bridge [OPEN]\\
Nejsilnější: Kronecker congruence filter $\nu_2>0 \iff p\equiv1\pmod4$ within genus-11 cluster\\
Celková klasifikace selekce $p=137$: MC slabý / OPEN (unique arithmetic filter, incomplete physical closure)\\
Predikce $\alpha^{-1}$:\\
\quad Bare: 137 (z $V_{\mathrm{eff}}$, podmíněně)\\
\quad S residuem: $137+f(137)=137.036541$ [MC]\\
\quad Experiment: 137.035999\\
\quad Přesnost (residual offset): 99.9996\% (difference $5.42\times10^{-4}$)
\end{statusbox}

\ifdefined\INCLUDEMODE
\else
  \end{document}
\fi
